	\section{Ciclo de vida del expediente}
	\label{sec:maquina_estados}
	El elemento que ordena todo el flujo de trabajo es el estatus de la queja. Se modeló como una
	máquina de estados con cuatro valores y transiciones dirigidas, decisión que se tomó por dos
	razones: permite validar en el servidor qué operaciones son admisibles en cada momento ---la
	regla RN-Q08 depende por completo de esto--- y da al quejoso una línea de tiempo comprensible
	del avance de su expediente.
	
	\begin{itemize}
		\item \texttt{RECIBIDA}. Estado inicial de toda queja, sin importar su origen. Es el único
		estado en el que el quejoso puede editar su propia queja, porque es el único en el que la
		Defensoría no ha empezado a trabajar sobre ella.
	
		\item \texttt{EN\_VALIDACION}. El personal de recepción tomó la queja y está verificando su
		documentación y sus antecedentes. La edición por parte del quejoso queda cerrada a partir de
		aquí, para evitar que el expediente cambie mientras se lo evalúa.
	
		\item \texttt{RECHAZADA}. Estado terminal. La queja no reunió los requisitos de
		documentación; el rechazo exige motivos registrados y notificación al quejoso.
	
		\item \texttt{TURNADA}. Estado terminal desde la perspectiva de recepción. La queja fue
		canalizada a un área y a un defensor, y recibió un folio oficial de turnado.
	\end{itemize}
	
	Los estados anteriores describen principalmente la fase de recepción y turnado, no el ciclo
	completo del expediente. Tras el turnado pueden intervenir los módulos de Primer Contacto,
	Subdefensoría y las instancias responsables de resolución, con estados y actuaciones propios.
	Por ello, la máquina aquí presentada debe interpretarse como un subproceso; el ciclo de vida
	completo se documenta conforme a las reglas de los módulos posteriores.
	
	\section{Diseño de la arquitectura}
	\label{sec:diseno_arquitectura}
	La arquitectura de la plataforma aplica los principios expuestos en la
	Sección~\ref{sec:arquitectura_software}. Esta sección describe cómo se concretaron: qué
	componentes existen, cómo se comunican, cómo se despliegan y qué alternativas se descartaron en
	cada decisión.
	
	\subsection{Vista de componentes}
	El sistema se descompuso en microservicios de dominio y servicios auxiliares de inteligencia
	artificial. Los servicios de negocio implementan las operaciones institucionales; los de IA
	atienden recuperación de información y redacción asistida. La
	Tabla~\ref{tab:microservicios} presenta los componentes principales. Los puertos de los
	servicios internos corresponden a la publicación en el servidor cuando aplica; el
	procesamiento inteligente se comunica por la red de contenedores y no requiere exposición pública.
	
	\rowcolors{2}{rowblue}{white}
	
	\small
	\renewcommand{\arraystretch}{1.25}
	\setlength{\tabcolsep}{4pt}
	
	\begin{longtable}{|p{3.4cm}|p{1.2cm}|p{6.6cm}|p{2.4cm}|}
		\caption{Microservicios de la plataforma y su responsabilidad}
		\label{tab:microservicios} \\
		\hline
		\rowcolor{headerblue}
		\headercell{Servicio} & \headercell{Puerto} & \headercell{Responsabilidad} &
		\headercell{Emite tokens} \\
		\hline
		\endfirsthead
	
		\hline
		\rowcolor{headerblue}
		\headercell{Servicio} & \headercell{Puerto} & \headercell{Responsabilidad} &
		\headercell{Emite tokens} \\
		\hline
		\endhead
	
		\texttt{auth-service} & 8083 & Sesión del quejoso, activación de cuenta, recuperación de
		contraseña y perfil & Sí, para quejosos \\ \hline
	
		\texttt{queja-service} & 8084 & Registro de quejas público y autenticado, consulta y edición
		de quejas propias, evidencias y respuesta a acuerdos de conciliación & No, solo verifica
		\\ \hline
	
		\texttt{notificaciones-\allowbreak service} & 8085 & Envío de correo electrónico y centro de
		notificaciones persistido por usuario & No, solo verifica \\ \hline
	
		\texttt{catalogo-\allowbreak service} & 8086 & Catálogo de dependencias y unidades académicas del
		Instituto & No, solo verifica \\ \hline
	
		\texttt{admin-service} & 8087 & Sesión del personal, cuentas y roles, plantillas de oficios,
		respaldos y bitácora de seguridad & Sí, para personal \\ \hline
	
		\texttt{revision-\allowbreak service} & 8088 & Bandeja de recepción, validación, rechazo, turnado,
		registro manual, historial exportable y emisión de acuerdos & No, solo verifica \\ \hline
	
		\texttt{chatbot-\allowbreak service} & 8089 & Menú de preguntas frecuentes del portal público y su
		administración & No, solo verifica \\ \hline
		\texttt{primer-contacto-service} & 8082 & Gestión de expedientes de Primer Contacto e integración
		de las herramientas de antecedentes y resúmenes & No, verifica \\ \hline
		\texttt{subdefensoria-service} & Interno & Operaciones del módulo de Subdefensoría & No, verifica \\ \hline
		\texttt{historico-service} & Interno & Consulta y tratamiento de información histórica
		de quejas & No, verifica \\ \hline
		\texttt{antecedentes-service} & Interno & Recuperación híbrida de expedientes con TF-IDF,
		MiniLM y fusión de rankings & No emite \\ \hline
		\texttt{resumen-service} & Interno & Selección extractiva y edición institucional de resúmenes &
		No emite \\ \hline
		\texttt{qwen-server} & Interno & Motor de inferencia del modelo generativo para el editor &
		No emite \\ \hline
	\end{longtable}
	
	\normalsize
	
	Los tokens de sesión se emiten desde los servicios de autenticación de quejosos y de
	personal. Los servicios de negocio que necesitan autorización los verifican; los componentes
	de IA permanecen accesibles únicamente a los servicios internos autorizados por la topología
	de red y las reglas de acceso. Esta separación impide confundir un servicio de inferencia
	con un punto de autenticación o una decisión administrativa.
	
	\subsection{Vista de despliegue}
	El despliegue diferencia el acceso público, los servicios de negocio y los componentes
	de procesamiento de inteligencia artificial. El entorno de aplicaciones aloja los
	microservicios y la persistencia operativa; el entorno de acceso sirve las aplicaciones
	web y el proxy inverso; el procesamiento de modelos dispone de servicios de inferencia
	separados y redes internas de contenedores. La Figura de despliegue y los detalles de
	infraestructura del capítulo de implementación especifican la distribución física y la
	configuración finalmente utilizada.
	
	Esta separación cumple tres propósitos. El primero es de seguridad: el servidor que contiene la
	base de datos y la lógica de negocio no está expuesto directamente a Internet, y su cortafuegos
	acepta conexiones a los puertos de los microservicios únicamente desde la dirección del servidor
	de acceso público. El segundo es de aislamiento de fallos: una saturación en la entrega de
	archivos estáticos no compite por recursos con el procesamiento de las peticiones de negocio.
	El tercero es operativo: el certificado de transporte y la configuración del dominio residen en
	un solo lugar, de modo que renovarlos no implica reiniciar ningún servicio de negocio.
	
	El único puerto realmente público es el de HTTPS; el puerto de HTTP se conserva abierto
	exclusivamente para redirigir el tráfico a su equivalente cifrado y para atender el reto de
	validación del certificado. La Tabla~\ref{tab:enrutamiento} presenta las reglas de enrutamiento
	del proxy inverso, que son el contrato entre el exterior y la arquitectura interna.
	
	\small
	\renewcommand{\arraystretch}{1.2}
	
	\begin{longtable}{|p{4.0cm}|p{5.4cm}|p{4.4cm}|}
		\caption{Reglas de enrutamiento del proxy inverso público}
		\label{tab:enrutamiento} \\
		\hline
		\rowcolor{headerblue}
		\headercell{Ruta pública} & \headercell{Destino} & \headercell{Servidor} \\
		\hline
		\endfirsthead
	
		\hline
		\rowcolor{headerblue}
		\headercell{Ruta pública} & \headercell{Destino} & \headercell{Servidor} \\
		\hline
		\endhead
	
		\texttt{/} & Portal del quejoso & Acceso público \\ \hline
		\texttt{/admin/} & Consola de administración & Acceso público \\ \hline
		\texttt{/revision/} & Panel de revisión & Acceso público \\ \hline
		\texttt{/primer-contacto/} & Panel de Primer Contacto & Acceso público \\ \hline
		\texttt{/subdefensoria/} & Panel de Subdefensoría & Acceso público \\ \hline
		\texttt{/api/auth/} & \texttt{auth-service} & Aplicación \\ \hline
		\texttt{/api/quejoso/} & \texttt{queja-service} & Aplicación \\ \hline
		\texttt{/api/notificaciones/} & \texttt{notificaciones-\allowbreak service} & Aplicación \\ \hline
		\texttt{/api/catalogos/} & \texttt{catalogo-\allowbreak service} & Aplicación \\ \hline
		\texttt{/api/admin/} & \texttt{admin-service} & Aplicación \\ \hline
		\texttt{/api/revision/} & \texttt{revision-\allowbreak service} & Aplicación \\ \hline
		\texttt{/api/chatbot/} & \texttt{chatbot-\allowbreak service} & Aplicación \\ \hline
		\texttt{/api/primer-contacto/} & \texttt{primer-contacto-service} & Aplicación \\ \hline
		\texttt{/api/subdefensoria/} & \texttt{subdefensoria-service} & Aplicación \\ \hline
	\end{longtable}
	
	\normalsize
	\setlength{\tabcolsep}{6pt}
	\renewcommand{\arraystretch}{1}
	
	Una consecuencia de diseño de este esquema merece destacarse: como el proxy entrega tanto las
	aplicaciones web como la interfaz de programación bajo un mismo dominio, el navegador considera
	que todo proviene del mismo origen y no activa la verificación de intercambio de recursos entre
	orígenes. Eso simplifica el modelo de seguridad del cliente sin sacrificar la separación en
	servicios del lado del servidor.
	
	\subsection{Diseño de la comunicación entre servicios}
	La comunicación entre servicios es síncrona sobre HTTP siguiendo el estilo REST, conforme a lo
	expuesto en la Sección~\ref{sec:arquitectura_software}. Se descartó la mensajería asíncrona
	porque el volumen de operaciones y el número de servicios no justifican la complejidad operativa
	de un intermediario de mensajes, y porque ninguno de los flujos identificados tolera bien la
	consistencia eventual: un quejoso que envía una queja espera su folio en la misma respuesta, no
	después.
	
	El diseño utiliza llamadas internas para validación de identidad y envío de notificaciones,
	además de consultas a los servicios especializados. En Primer Contacto, la recuperación de
	antecedentes se obtiene mediante una petición HTTP al servicio de búsqueda; la elaboración de
	resúmenes invoca otro servicio, que a su vez consulta un motor de inferencia generativa.
	Cada interacción tiene su propio tiempo de espera y una estrategia de manejo de fallos: la
	indisponibilidad del componente de IA no debe sustituir ni bloquear el criterio institucional.
	Las dependencias de datos y sus compromisos se analizan en la
	Sección~\ref{sec:decisiones_diseno}.
	
	\subsection{Diseño de la seguridad y el control de acceso}
	El esquema de seguridad se construyó sobre cuatro mecanismos, cada uno atendiendo una amenaza
	distinta.
	
	La \textbf{autenticación} emplea tokens firmados con algoritmo simétrico. Un token emitido por
	el servicio de autenticación es válido en los servicios que atienden al quejoso, y uno emitido
	por el servicio de administración es válido en los que atienden al personal, porque todos
	comparten la misma clave de firma. Cada servicio verifica el token de forma local, sin consultar
	al emisor, lo que evita un punto único de falla en cada petición. La vigencia del token es de
	una hora, y el cliente cierra la sesión de forma automática cuando el servidor la rechaza.
	
	El \textbf{almacenamiento de contraseñas} usa una función de resumen con costo configurable en
	los dos servicios que emiten sesión. La política de complejidad exige un mínimo de ocho
	caracteres, al menos una mayúscula y al menos un dígito, y está implementada como una expresión
	compartida para que no divergja entre servicios. Las cuentas de personal se crean con contraseña
	temporal y obligación de cambiarla en el primer acceso.
	
	La \textbf{autorización} se declara por rol en cada controlador, no solo por ruta. Esta
	distinción no es cosmética: una restricción declarada únicamente en la configuración de rutas
	puede eludirse si una ruta nueva se agrega sin actualizar esa configuración, mientras que la
	declaración en el controlador viaja con el código que protege.
	
	El \textbf{cifrado en tránsito} se termina en el proxy inverso con un certificado de una
	autoridad reconocida y renovación automática. Del proxy hacia los microservicios el tráfico
	viaja sin cifrar, decisión aceptable porque ese tramo no atraviesa Internet y está restringido
	por dirección de origen en el cortafuegos, pero que conviene revisar si en el futuro los
	servidores dejan de estar bajo control directo del proyecto.
	
	Un punto pendiente que debe declararse: el cierre automático de sesión por inactividad
	solicitado para las cuentas de personal no está implementado al cierre de esta etapa. La
	expiración del token limita la ventana de exposición a una hora, pero no equivale a un cierre
	por inactividad.
	
	\subsection{Decisiones de diseño y alternativas descartadas}
	\label{sec:decisiones_diseno}
	Esta sección documenta las decisiones estructurales del proyecto con el formato
	contexto-alternativas-decisión-consecuencias. Registrarlas de esta forma cumple un propósito
	concreto: una decisión de arquitectura cuyo razonamiento no se documenta se vuelve
	indistinguible de un accidente, y quien mantenga el sistema en el futuro no sabrá si puede
	revertirla.
	
	\paragraph{Decisión 1: base de datos única compartida.}
	\emph{Contexto:} varios servicios de negocio necesitan datos que se relacionan entre sí ---una
	queja, su tutor, sus evidencias, el usuario que la presentó, la notificación que se le envió---.
	\emph{Alternativas:} una base de datos por servicio, con sincronización explícita entre ellas;
	o una base única compartida. \emph{Decisión:} base única. \emph{Consecuencias aceptadas:} se
	sacrifica la autonomía de datos que caracteriza a los microservicios puros, y dos servicios
	escriben la misma tabla sin coordinación transaccional entre ellos. A cambio se evita la
	complejidad de mantener consistencia entre bases separadas, que a la escala de este sistema
	habría requerido un esfuerzo mayor que el beneficio obtenido. Es un híbrido consciente, no una
	implementación incompleta del patrón.
	
	\paragraph{Decisión 2: entidades independientes sobre una tabla compartida.}
	\emph{Contexto:} el servicio de quejas y el de revisión necesitan la misma tabla de quejas, pero
	con vistas distintas: el primero requiere los campos del quejoso, el segundo los del flujo de
	validación. \emph{Alternativas:} un servicio dueño de la tabla que exponga los datos al otro por
	interfaz de programación; o dos clases de persistencia independientes sobre la misma tabla
	física. \emph{Decisión:} entidades independientes. \emph{Consecuencias aceptadas:} un cambio de
	estatus hecho por recepción es visible de inmediato para el quejoso sin llamada entre servicios,
	lo que simplifica el flujo; pero ambos servicios pueden agregar columnas sin que el otro lo
	sepa, y una modificación descoordinada del esquema puede romper al servicio vecino. La mitigación
	adoptada es que cada servicio declare únicamente las columnas que usa.
	
	\paragraph{Decisión 3: catálogo institucional como servicio propio.}
	\emph{Contexto:} el catálogo de dependencias del Instituto es requerido por el formulario público
	de quejas, por el turnado en el panel de revisión y por la consola de administración.
	\emph{Alternativas:} alojarlo dentro del servicio de quejas, que fue quien primero lo necesitó;
	o extraerlo a un servicio propio. \emph{Decisión:} servicio propio.
	\emph{Consecuencias aceptadas:} se añadió un servicio más que operar y desplegar, con el
	sobrecosto correspondiente, en un momento en que un solo consumidor lo necesitaba. La razón para
	asumirlo fue anticipada: al convertirse en dato de referencia de tres consumidores distintos,
	alojarlo dentro de un servicio de negocio habría obligado a que los otros dos dependieran de él
	por razones ajenas a su dominio.
	
	\paragraph{Decisión 4: evidencias almacenadas en la base de datos.}
	\emph{Contexto:} cada queja puede llevar varios archivos de evidencia que constituyen material
	probatorio. \emph{Alternativas:} guardar los archivos en el sistema de archivos del servidor y
	almacenar la ruta en la base; o almacenar el contenido binario completo dentro de la base.
	\emph{Decisión:} contenido binario en la base. \emph{Consecuencias aceptadas:} la base crece más
	rápido y los respaldos son más grandes. A cambio se obtiene lo que en este dominio pesa más: el
	archivo y su expediente se respaldan y se restauran como una unidad atómica, de modo que no puede
	existir un expediente cuya evidencia se haya perdido porque el respaldo del sistema de archivos y
	el de la base se tomaron en momentos distintos. Para material con posible valor probatorio, esa
	garantía justifica el costo de almacenamiento.
	
	\paragraph{Decisión 5: aplicaciones web separadas por perfil funcional.}
	\emph{Contexto:} los módulos de quejosos, recepción, administración, Primer Contacto y
	Subdefensoría tienen responsabilidades y necesidades de interfaz distintas.
	\emph{Alternativas:} una sola aplicación con rutas condicionadas por rol o aplicaciones
	separadas. \emph{Decisión:} separar los módulos web por ámbito de uso.
	\emph{Consecuencias aceptadas:} se deben mantener coherentes componentes, interceptores
	y estilos entre proyectos; a cambio se delimitan interfaces y despliegues de cada área.
	La autorización efectiva sigue dependiendo de la validación del servidor.
	
	\paragraph{Decisión 6: datos del quejoso en columnas propias.}
	\emph{Contexto:} el registro público debe capturar la identidad completa del quejoso, que en una
	versión inicial se concatenaba como texto libre dentro de la descripción de la queja.
	\emph{Alternativas:} conservar la concatenación, más simple de implementar; o estructurar cada
	dato en su propia columna. \emph{Decisión:} columnas propias. \emph{Consecuencias aceptadas:} el
	esquema creció y el formulario requiere más validación. A cambio, los datos se volvieron
	consultables y verificables ---la búsqueda de antecedentes del requerimiento RF-R04 sería
	impracticable sobre texto concatenado--- y la unidad académica pudo pasar de texto libre a una
	referencia al catálogo institucional.
	
	\paragraph{Decisión 7: una imagen de contenedor por servicio.}
	\emph{Contexto:} los microservicios comparten la misma imagen base y el mismo procedimiento de
	construcción. \emph{Alternativas:} una imagen común reutilizada por todos, parametrizada al
	ejecutarse; o una imagen dedicada por servicio. \emph{Decisión:} una por servicio.
	\emph{Consecuencias aceptadas:} se consume más espacio de almacenamiento en el registro de
	imágenes. A cambio, cada contenedor es identificable y versionable de forma independiente, y se
	elimina la ambigüedad de no poder saber qué versión de qué servicio está ejecutando un
	contenedor determinado.
	
	\section{Diseño de la base de datos}
	\label{sec:diseno_bd}
	La persistencia principal de expedientes es relacional y se comparte entre varios servicios
	de negocio. En algunos casos, distintas entidades de persistencia trabajan sobre la misma
	tabla física, bajo el patrón descrito en la Decisión~2. La
	Tabla~\ref{tab:esquema_bd} presenta las tablas de los módulos de registro, revisión y
	administración. Las estructuras adicionales de Primer Contacto, Subdefensoría e histórico
	deben documentarse en el diccionario de datos correspondiente. Los embeddings del recuperador
	se conservan en una caché técnica independiente, que no reemplaza el expediente original.
	
	\rowcolors{2}{rowblue}{white}
	
	\small
	\renewcommand{\arraystretch}{1.25}
	\setlength{\tabcolsep}{4pt}
	
	\begin{longtable}{|p{3.5cm}|p{3.1cm}|p{7.5cm}|}
		\caption{Esquema de la base de datos por tabla y servicio propietario}
		\label{tab:esquema_bd} \\
		\hline
		\rowcolor{headerblue}
		\headercell{Tabla} & \headercell{Propietario} & \headercell{Contenido} \\
		\hline
		\endfirsthead
	
		\hline
		\rowcolor{headerblue}
		\headercell{Tabla} & \headercell{Propietario} & \headercell{Contenido} \\
		\hline
		\endhead
	
		\texttt{usuarios} & \texttt{auth-service} & Cuentas de quejoso: identidad institucional,
		credencial cifrada, datos de contacto, datos de tutor y código de recuperación con su
		vigencia \\ \hline
	
		\texttt{quejas} & \texttt{queja-service} y \texttt{revision-\allowbreak service} & Cuerpo del expediente:
		folio, identidad del quejoso, narrativa de los hechos, unidad académica, denunciado, origen,
		estatus, y los campos del flujo de validación ---motivo de rechazo, área, defensor asignado,
		quién validó y cuándo--- \\ \hline
	
		\texttt{queja\_tutores} & \texttt{queja-service} & Datos del tutor o adulto responsable,
		en relación uno a uno con la queja \\ \hline
	
		\texttt{queja\_evidencias} & \texttt{queja-service} & Archivos de evidencia: nombre, tipo,
		tamaño y contenido binario completo \\ \hline
	
		\texttt{acuerdos\_} \texttt{conciliacion} & \texttt{revision-\allowbreak service} y
		\texttt{queja-service} & Acuerdos emitidos por la Defensoría: asunto, términos, estado y
		respuesta del quejoso \\ \hline
	
		\texttt{notificaciones} & \texttt{notificaciones-} \texttt{service} & Avisos persistidos por
		destinatario, con tipo, estado de lectura y enlace \\ \hline
	
		\texttt{dependencias} & \texttt{catalogo-\allowbreak service} & Catálogo institucional jerárquico con
		clave propia, clave del padre, tipo, categoría y nivel \\ \hline
	
		\texttt{personal\_} \texttt{administrativo} & \texttt{admin-service} & Cuentas de personal:
		identidad, rol, credencial cifrada, condición de cuenta temporal y último acceso \\ \hline
	
		\texttt{plantillas\_} \texttt{documentos} & \texttt{admin-service} & Plantillas de oficios
		con clave estable de negocio, contenido y datos de última actualización \\ \hline
	
		\texttt{bitacora\_acciones} & \texttt{admin-service} & Registro de auditoría: usuario, acción,
		dirección de origen y fecha \\ \hline
	
		\texttt{preguntas\_chatbot} & \texttt{chatbot-\allowbreak service} & Contenido del menú de preguntas
		frecuentes del portal público \\ \hline
	\end{longtable}
	
	\normalsize
	
	\subsection{Correspondencia entre reglas de negocio y mecanismos de integridad}
	Un modelo de datos que solo almacena información deja las reglas del dominio a merced de que
	cada consumidor las respete. El diseño de esta base busca lo contrario: que las reglas se hagan
	cumplir en el nivel más bajo posible, de modo que ningún camino de acceso pueda violarlas. La
	Tabla~\ref{tab:reglas_mecanismos} muestra esa correspondencia para las reglas que admiten
	verificación estructural.
	
	\small
	\renewcommand{\arraystretch}{1.2}
	
	\begin{longtable}{|p{1.9cm}|p{5.6cm}|p{6.6cm}|}
		\caption{Reglas de negocio y mecanismo que las hace cumplir}
		\label{tab:reglas_mecanismos} \\
		\hline
		\rowcolor{headerblue}
		\headercell{Regla} & \headercell{Enunciado} & \headercell{Mecanismo} \\
		\hline
		\endfirsthead
	
		\hline
		\rowcolor{headerblue}
		\headercell{Regla} & \headercell{Enunciado} & \headercell{Mecanismo} \\
		\hline
		\endhead
	
		RN-Q04 & El par folio-correo es la llave de activación & Consulta entre servicios que exige
		coincidencia de ambos campos antes de crear la cuenta \\ \hline
	
		RN-Q05 & Una cuenta activa no se reactiva & Verificación del indicador de actividad antes de
		procesar la activación \\ \hline
	
		RN-Q07 & El código de recuperación expira en diez minutos & Columna de fecha de expiración
		evaluada en el servidor al validar el código \\ \hline
	
		RN-Q08 & Una queja solo es editable en estatus recibida & Validación del estatus en la capa
		de negocio antes de aplicar la modificación \\ \hline
	
		RN-Q09 & Solo el destinatario responde un acuerdo & Comparación del correo del acuerdo contra
		el del token de sesión \\ \hline
	
		RN-Q10 & El origen de la queja se infiere & Asignación en el servidor según el punto de acceso
		invocado, sin exponerlo como campo de entrada \\ \hline
	
		RN-R01 & Transiciones válidas del expediente & Columna de estatus con conjunto acotado de
		valores y validación de la transición en la capa de negocio \\ \hline
	
		RN-R02 & Solo recepción usa el servicio de revisión & Declaración de rol requerido en cada
		controlador \\ \hline
	
		RN-R06 & Solo defensores activos reciben turnado & Consulta filtrada por rol y por indicador
		de actividad \\ \hline
	
		RN-A03 & La baja desactiva, no elimina & Indicador de actividad en lugar de borrado físico,
		preservando la integridad referencial con la bitácora \\ \hline
	
		RN-A04 & Correo y número de empleado únicos & Restricción de unicidad en la definición de la
		tabla \\ \hline
	
		RN-A06 & La plantilla se identifica por clave de negocio & Restricción de unicidad sobre la
		columna de tipo, independiente del identificador autogenerado \\ \hline
	\end{longtable}
	
	\normalsize
	\setlength{\tabcolsep}{6pt}
	\renewcommand{\arraystretch}{1}
	
	El diccionario de datos completo, con el tipo y la longitud de cada columna, las cardinalidades
	detalladas, los índices y las sentencias de creación, se documenta en el anexo correspondiente.
	
	\section{Diseño de los componentes de inteligencia artificial}
	\label{sec:diseno_ia}
	Los componentes inteligentes se diseñaron como herramientas auxiliares para la revisión de
	expedientes. Se separaron en dos procesos con objetivos y salidas diferentes: recuperación
	de antecedentes relacionados y elaboración de borradores narrativos. El primero devuelve
	referencias documentales ordenadas por relevancia; el segundo devuelve un texto redactado
	a partir de información seleccionada de la queja. Ninguno determina automáticamente la
	competencia institucional, la existencia de una duplicidad ni el sentido de una resolución.

	\subsection{Diseño del componente de búsqueda de antecedentes}
	\label{sec:diseno_antecedentes}
	La entrada principal es la narración de una queja. El diseño parte de una colección de
	quejas vigentes e históricas que se representa con dos señales complementarias:
	coincidencias de términos y similitud semántica. La primera emplea TF-IDF con unigramas y
	bigramas; la segunda emplea un codificador multilingüe MiniLM para construir embeddings
	de fragmentos del contenido.

	\paragraph{Fragmentación de narrativas.}
	Las narrativas extensas se dividen en segmentos de longitud acotada, con solapamiento entre
	segmentos consecutivos. El objetivo es preservar contexto local y evitar que un solo vector
	de documento oculte acontecimientos relevantes mencionados en una parte de la queja.
	Las decisiones de tamaño, solapamiento y procesamiento por lotes pertenecen a la
	configuración de implementación y se evalúan según costo y calidad de recuperación.

	\paragraph{Recuperación paralela y fusión de rankings.}
	Cada recuperador produce su propia lista ordenada de antecedentes: TF-IDF privilegia
	coincidencias léxicas específicas, mientras que MiniLM facilita localizar textos de
	significado relacionado aunque utilicen otras expresiones. Se emplea Reciprocal Rank
	Fusion (RRF) ponderado, que integra las posiciones en ambos rankings en lugar de sumar
	directamente sus puntuaciones, porque las escalas de similitud de los recuperadores no
	son necesariamente equivalentes. El mecanismo se fundamenta en la
	Sección~\ref{sec:busqueda_antecedentes}.

	La configuración seleccionada para la integración del sistema asigna un peso
	del 75\,\% al ranking TF-IDF y del 25\,\% al ranking semántico MiniLM.
	Estas ponderaciones se aplican a las contribuciones de los rankings mediante RRF;
	no representan porcentajes de similitud ni probabilidades de duplicidad.
	Su selección corresponde a los experimentos con datos sintéticos documentados en
	el capítulo de pruebas, y no implica que constituyan una combinación óptima para
	todos los escenarios institucionales. Para compararlos es necesario utilizar un mismo
	conjunto de consultas y juicios de relevancia, probar distintas ponderaciones y medir
	la pertinencia y posición de los antecedentes mediante métricas como precisión,
	exhaustividad y ganancia acumulada descontada normalizada.

	\paragraph{Indexación y actualización.}
	El diseño evita recalcular los embeddings de toda la colección en cada actualización:
	mantiene una caché persistente de representaciones, identificadas mediante una huella
	del contenido y la configuración del codificador. Si un documento no cambió,
	se puede reutilizar su representación; si cambió o se incorporó, se calcula la nueva.
	El índice disponible para consulta se sustituye cuando la reconstrucción termina,
	para no exponer un estado parcialmente construido.

	\paragraph{Presentación e interpretación.}
	Los antecedentes se exponen con su identificación y datos contextuales disponibles,
	para que el personal contraste los hechos. La clasificación de coincidencias por hechos
	o involucrados sirve como apoyo visual. Una similitud elevada no confirma que dos
	quejas sean idénticas: la identificación de duplicidades corresponde a la revisión humana.

	\subsection{Diseño del componente de generación de resúmenes}
	\label{sec:diseno_resumenes}
	La generación de resúmenes sigue un esquema de dos etapas para separar la selección de
	información de la redacción. La entrada es la narrativa original de la queja y la salida
	es un borrador institucional, nunca una sustitución del documento fuente.

	\paragraph{Selección extractiva.}
	En la primera etapa, el procedimiento de selección extractiva utiliza criterios de relevancia,
	redundancia, importancia y conservación del contexto. Incorpora un criterio de selección
	basado en Maximum Marginal Relevance (MMR), de modo que el resultado no concentre
	exclusivamente oraciones muy similares. Su salida es una captura extractiva, que
	conserva fragmentos del relato en su formulación original y sirve de entrada al editor.

	\paragraph{Edición mediante un modelo de lenguaje.}
	La segunda etapa utiliza Qwen3, un modelo generativo, para convertir la captura
	extractiva en narración institucional en tercera persona. La instrucción de edición
	busca conservar actores, negaciones, fechas, cantidades, evidencias, consecuencias y
	secuencia de acontecimientos, así como presentar las afirmaciones como declaraciones
	del quejoso y no como hechos jurídicamente comprobados. Su función prevista es
	reformular los fragmentos recibidos, no realizar una nueva selección deliberada
	de hechos. Sin embargo, debido al carácter generativo del modelo, pueden ocurrir
	omisiones o alteraciones de información incluso bajo estas instrucciones.

	\paragraph{Procesamiento por bloques y límites operativos.}
	La captura puede dividirse en bloques de oraciones con una extensión máxima, para
	acotar el trabajo por solicitud al motor generativo y adecuar el procesamiento a la
	infraestructura de CPU disponible. Se establece un presupuesto de tiempo para la
	edición y se prevé conservar el fragmento extractivo si una llamada falla o queda
	incompleta. Esa contingencia evita descartar contenido por errores de generación,
	pero puede producir un borrador parcialmente editado, por lo que debe informarse
	en las pruebas y en la interpretación de resultados.

	\paragraph{Fidelidad y validación humana.}
	El diseño distingue dos fuentes posibles de omisión: la selección extractiva y la
	edición generativa. Es necesario evaluar ambas por separado, comparando la narrativa
	original con la captura seleccionada y esta con el resumen final. Se consideran
	críticos los cambios de actores, negaciones, fechas y cantidades. Aunque las
	instrucciones de generación buscan evitar alteraciones, no ofrecen una garantía
	de conservación factual; toda síntesis debe quedar sujeta a revisión humana.
	Por restricciones de tiempo, la evaluación realizada con narrativas sintéticas
	no se complementó con una validación de calidad y utilidad mediante expedientes
	reales y la participación del personal de la Defensoría. Por tanto, las pruebas
	técnicas no deben interpretarse como validación institucional del componente.
	La revisión de calidad con datos autorizados y personal responsable queda como
	actividad futura.

	\subsection{Integración de los servicios con Primer Contacto}
	\label{sec:integracion_ia}
	El módulo de Primer Contacto actúa como consumidor de ambos componentes. El servicio
	de antecedentes atiende peticiones de consulta y devuelve expedientes sugeridos;
	el servicio de resúmenes ejecuta la selección extractiva y solicita la redacción
	al motor generativo. La separación de los servicios permite actualizar los
	algoritmos de recuperación y generación sin acoplarlos entre sí.

	Por tratarse de llamadas síncronas, el diseño contempla tiempos máximos de espera,
	observabilidad del estado de los servicios y manejo de fallos. Si un componente no
	responde, la actividad administrativa debe poder continuar sin interpretar el
	resultado faltante como una decisión automatizada. La descripción de contratos
	HTTP, parámetros y despliegue corresponde al capítulo de implementación.

	\subsection{Decisiones de diseño y alternativas consideradas}
	La Tabla~\ref{tab:decisiones_ia} resume los principales compromisos del diseño.
	\small
	\renewcommand{\arraystretch}{1.2}
	\begin{longtable}{|p{3.1cm}|p{4.5cm}|p{6.1cm}|}
		\caption{Decisiones técnicas de los componentes de inteligencia artificial}
		\label{tab:decisiones_ia} \\
		\hline
		\rowcolor{headerblue}
		\headercell{Decisión} & \headercell{Alternativas} & \headercell{Justificación y compromiso} \\ \hline
		\endfirsthead
		\hline
		\rowcolor{headerblue}
		\headercell{Decisión} & \headercell{Alternativas} & \headercell{Justificación y compromiso} \\ \hline
		\endhead
		Recuperación híbrida & TF-IDF aislado, MiniLM aislado o combinación &
		Combina coincidencias específicas y relaciones semánticas; exige evaluar sus pesos. \\ \hline
		Fragmentación & Vector único por documento o vectores por fragmento &
		Conserva señales locales de textos extensos; aumenta almacenamiento y cómputo. \\ \hline
		Fusión por rankings & Suma directa de puntuaciones o RRF ponderado &
		Evita comparar directamente escalas de similitud no equivalentes; se seleccionó
		la ponderación 75\,\% TF-IDF y 25\,\% MiniLM con evidencia sintética,
		sin asumir generalización a expedientes reales. \\ \hline
		Caché de embeddings & Reconstrucción total o reutilización por contenido &
		Reduce trabajo repetido; obliga a invalidar representaciones al cambiar contenido o modelo. \\ \hline
		Resumen en dos etapas & Generación directa o selección más edición &
		Separa relevancia y redacción; conserva riesgos de omisión en ambas etapas. \\ \hline
		Inferencia en CPU & Modelo de mayor tamaño o modelo compacto cuantizado &
		Permite operación en los recursos disponibles; limita concurrencia y velocidad. \\ \hline
	\end{longtable}
	\normalsize
	\renewcommand{\arraystretch}{1}

	\section{Diseño de la interfaz de usuario}
	\label{sec:diseno_interfaz}
	El diseño de las interfaces partió de prototipos de alta fidelidad que se usaron como
	herramienta de validación temprana con el cliente. La razón de invertir en ellos antes de
	escribir código fue económica: corregir el orden de los campos de un formulario en un prototipo
	cuesta minutos, y hacerlo sobre una interfaz ya implementada y desplegada cuesta un ciclo
	completo de compilación y despliegue.
	
	Esa inversión se justificó en la práctica. La retroalimentación de las sesiones de revisión
	derivó en dos tipos de cambio que se incorporaron antes de implementar: ajustes en el orden y
	agrupación de los campos del formulario de queja, y la incorporación de mensajes de validación
	explícitos donde el prototipo original solo mostraba el campo en rojo.
	
	Los criterios de diseño aplicados a las aplicaciones web fueron los siguientes:
	
	\begin{itemize}
		\item \textbf{Identidad institucional.} La paleta y la estructura del encabezado y del pie
		siguen la identidad visual del Instituto, de modo que el usuario reconozca que se encuentra en
		un servicio oficial y no en un sitio de terceros.
	
		\item \textbf{Formularios por pasos.} Los formularios extensos se dividen en secciones
		numeradas en lugar de presentarse como una lista continua de campos, para reducir la carga
		cognitiva y hacer visible el progreso.
	
		\item \textbf{Ayuda contextual no intrusiva.} Los avisos normativos que el usuario debe
		conocer ---como el plazo para presentar una queja--- se presentan en un elemento flotante
		consultable en cualquier momento, en lugar de un bloque fijo que reduce el espacio útil del
		formulario.
	
		\item \textbf{Retroalimentación explícita.} Toda operación produce un mensaje visible de
		éxito, error o validación pendiente. Este criterio se elevó a requerimiento no funcional
		(RNF-Q07) porque su ausencia fue la causa de que, en una versión temprana, los mensajes de
		error existieran en el código pero no fueran visibles para el usuario.
	
		\item \textbf{Consistencia entre paneles de personal.} Las interfaces de administración,
		revisión, Primer Contacto y Subdefensoría comparten criterios visuales, patrones de
		error y estilos institucionales, para reducir diferencias innecesarias entre áreas.
		\item \textbf{Resultados de inteligencia artificial revisables.} La interfaz separa los
		antecedentes relacionados del borrador de resumen, identifica sus resultados como apoyos
		para el análisis y permite consultar la narrativa original al revisar el contenido.
	\end{itemize}
	
	Los prototipos completos y las capturas de las interfaces implementadas se presentan en el
	Capítulo~\ref{cap:modulos} y en el anexo correspondiente.
